\newpage
\section{Wstęp}

Pandemia wywołana przez wirusa SARS-CoV-2, ogłoszona przez Światową Organizację Zdrowia w~marcu 2020~roku, stanowiła jedno z~najpoważniejszych zagrożeń zdrowia publicznego w~najnowszej historii. Równocześnie sezonowe epidemie grypy, wywoływane między innymi przez wirusy Influenza~A, od dziesięcioleci pozostają istotnym problemem epidemiologicznym, powodując setki tysięcy zgonów rocznie na całym świecie. W~obu przypadkach kluczową rolę w~procesie infekcji odgrywają białka powierzchniowe wirusa --- białko Spike~(S) w~przypadku SARS-CoV-2 oraz hemaglutynina~(HA) w~przypadku wirusa grypy~\cite{walls2020spike, bush1999influenza}. Białka te są głównym celem odpowiedzi immunologicznej gospodarza, a~zarazem stanowią podstawowy antygen wykorzystywany w~konstrukcji szczepionek. Mutacje zachodzące w~regionach epitopowych tych białek prowadzą do zjawiska ucieczki immunologicznej (ang.~\textit{immune escape}), w~wyniku którego wcześniej nabyta odporność --- zarówno poszczepienną, jak i~poinfekcyjna --- może okazać się niewystarczająca wobec nowych wariantów wirusa~\cite{harvey2021variants}. Konsekwencją tego procesu jest konieczność cyklicznej aktualizacji składu szczepionek, co wymaga ciągłego monitorowania ewolucji genomowej patogenów.

Dynamiczny rozwój technologii sekwencjonowania nowej generacji doprowadził do bezprecedensowego wzrostu liczby dostępnych danych sekwencyjnych. Baza GISAID (ang.~\textit{Global Initiative on Sharing All Influenza Data}) zgromadziła ponad 15~milionów sekwencji genomowych wirusa SARS-CoV-2, natomiast baza NCBI Influenza Virus Resource udostępnia wieloletnie dane dotyczące wirusów grypy~\cite{elbe2017gisaid, bao2008ncbi}. Tak ogromna skala danych sprawia, że tradycyjne metody analizy filogenetycznej stają się niewystarczające, co motywuje poszukiwanie podejść opartych na uczeniu maszynowym. W~ostatnich latach zaproponowano szereg metod predykcji ewolucji wirusów, obejmujących zarówno modele językowe białek~\cite{hie2021viral}, jak i~podejścia statystyczne identyfikujące czynniki napędzające pojawianie się nowych wariantów~\cite{maher2022predicting}. Niemniej jednak niewiele istniejących rozwiązań łączy reprezentacje wektorowe sekwencji białkowych (ang.~\textit{embeddings}) z~rekurencyjnymi sieciami neuronowymi wyposażonymi w~mechanizmy atencji, osadzonymi w~kontekście analizy szeregów czasowych.

Zdolność do przewidywania, w~których regionach białka wirusowego zajdą mutacje w~nadchodzących okresach, ma istotne znaczenie praktyczne. Umożliwia ona proaktywne projektowanie szczepionek, skracając czas reakcji na pojawienie się nowych wariantów. Ponadto istnieje zapotrzebowanie na uniwersalne narzędzia, które nie będą ograniczone do pojedynczego patogenu, lecz pozwolą na analizę różnych białek wirusowych przy minimalnych modyfikacjach konfiguracji. Podejście oparte na klasteryzacji sekwencji w~oknach czasowych, w~połączeniu ze śledzeniem ewolucji klastrów między kolejnymi okresami, oferuje naturalny sposób modelowania dynamiki zmian w~populacji wirusowej.

Wkład niniejszej pracy obejmuje pięć głównych elementów. Po pierwsze, zaprojektowano i~zaimplementowano uniwersalny framework do predykcji mutacji aminokwasowych w~białkach wirusowych, konfigurowalny za pomocą plików YAML dla różnych patogenów. Po drugie, opracowano wieloetapowy potok przetwarzania danych, transformujący surowe pliki FASTA w~gotowy zbiór danych treningowych poprzez czyszczenie sekwencji, podział na okresy czasowe, klasteryzację K-means z~wykorzystaniem reprezentacji ProtVec~\cite{asgari2015protvec} oraz dwukierunkowe łączenie klastrów między kolejnymi okresami. Po trzecie, zaimplementowano i~porównano trzy architektury modeli opartych na rekurencyjnych sieciach neuronowych: bazowy model LSTM~\cite{hochreiter1997lstm}, model LSTM z~mechanizmem atencji Bahdanau~\cite{bahdanau2015attention} oraz model LSTM z~podwójną atencją (ang.~\textit{Dual-Stage Attention-Based RNN}, DA-RNN)~\cite{qin2017darnn}. Po czwarte, przeprowadzono ewaluację eksperymentalną na dwóch zbiorach danych --- sekwencjach białka~S wirusa SARS-CoV-2 oraz białka HA wirusa Influenza~A podtypu H1N1 --- stosując chronologiczny podział danych oraz współczynnik korelacji Matthewsa (ang.~\textit{Matthews Correlation Coefficient}, MCC) jako główną metrykę oceny~\cite{chicco2020mcc}. Po piąte, zapewniono reprodukowalność eksperymentów poprzez konteneryzację środowiska z~wykorzystaniem platformy Docker.

Niniejsza praca składa się z~sześciu rozdziałów. Rozdział~2 formułuje cel pracy, tezę badawczą oraz precyzuje zakres przeprowadzonych badań. Rozdział~3 zawiera przegląd literatury obejmujący podstawy biologiczne, metody reprezentacji sekwencji białkowych, algorytmy klasteryzacji oraz architektury rekurencyjnych sieci neuronowych z~mechanizmami atencji. W~rozdziale~4 przedstawiono szczegółowy opis zaproponowanego rozwiązania, obejmujący architekturę frameworka, potok przetwarzania danych oraz implementację modeli. Rozdział~5 prezentuje wyniki ewaluacji eksperymentalnej na obu zbiorach danych wraz z~analizą porównawczą zastosowanych architektur. Rozdział~6 zawiera podsumowanie pracy, krytyczną analizę osiągniętych rezultatów oraz propozycje kierunków dalszych badań.
